\section{Summary of Findings}

This thesis investigated how a proactive, context-sensitive AI system can support goal-oriented remote
meetings during convergent decision-making processes. It addressed three research questions, moving from the
identification of facilitation strategies and challenges to the design and empirical evaluation of an
AI-based facilitation prototype.

\textbf{Regarding RQ1, facilitating strategies and challenges in goal-oriented remote meetings,} the literature review
and three expert interviews highlighted that effective decision-making requires more than a clearly defined
agenda. Facilitators must maintain focus, structure the decision-making process, ensure that relevant
information is considered, and create an environment in which participants can express differing
perspectives. Particular challenges arise from unequal participation, premature consensus, and the tendency
to discuss shared information more extensively than information held by individual participants. The
interviews also revealed considerable reservations about proactive AI interventions. Concerns included
inappropriate timing, the potential to embarrass participants, insufficient contextual understanding, and the
risk of undermining group autonomy. While AI was considered promising for structured and clearly defined
tasks, the experts emphasized that human facilitation involves social and contextual competencies that are
difficult to replicate technically.

Building on these findings, \textbf{RQ2, concerning the design requirements for an AI-based facilitation system} , was
addressed through the development of a scoped prototype. Rather than attempting to replicate the full range
of facilitation activities, the system focused on hidden-profile scenarios, in which relevant information is
distributed asymmetrically among participants. The prototype was designed to support balanced participation,
encourage the sharing of unique information, and preserve dissent to counteract premature consensus. Its
architecture adapted the Inner Thoughtsframework to enable the AI to continuously assess the ongoing
conversation, evaluate potential facilitation interventions, and decide whether to contribute proactively. At
the same time, the system was designed to remain neutral regarding the substantive decision and to avoid
imposing its own preferences. The development process demonstrated how established facilitation principles
can be translated into concrete system requirements and implemented in an LLM-based prototype. However, the
findings also underline that context-sensitive facilitation requires more than generating relevant questions:
the system must recognize when an intervention is needed and deliver it at an appropriate moment, which
remains challenging.

\textbf{RQ3 examined how participants perceived meetings involving a proactive AI moderator.}
The prototype was evaluated in three hidden-profile group sessions involving twelve participants. The results revealed a
discrepancy between the intended benefits of proactive facilitation and the actual interaction experience.
Participants generally perceived the quality of discussion positively, and no consistently negative pattern
emerged across the measured dimensions of communication, status effects, teamwork, and efficiency.
Nevertheless, the AI's interventions were predominantly perceived as disruptive. Participants considered the
system to intervene too frequently, often at inappropriate moments or with delays that made its contributions
less relevant to the ongoing discussion. Some interventions also demonstrated insufficient contextual
understanding, including misunderstandings of information about the candidates.
At the same time, participants recognized the potential value of the system's intended functions. Addressing
quieter participants and explicitly requesting additional perspectives were perceived as helpful in some
situations. The findings therefore suggest that the underlying facilitation objectives were not inherently
rejected; rather, their implementation frequently failed to provide value at the right moment. Although the
system addressed several of its intended user stories, only one of the three groups selected the optimal
candidate. This outcome cannot be attributed to the AI's interventions alone, but it illustrates that the
presence of a proactive moderator does not automatically ensure a well-founded group decision.

Taken together, the findings indicate that the central challenge of proactive AI facilitation is not simply
determining what to say, but deciding whether, when, and how to intervene. Proactive systems may offer value
where relevant information would otherwise remain unshared or where premature consensus threatens decision
quality. However, poorly timed or repetitive contributions can undermine these benefits by interrupting
communication and reducing participants' willingness to engage with the system. The results therefore support
the continued investigation of proactive AI facilitation while emphasizing the importance of contextual
accuracy, restrained intervention, and human control. The prototype represents an initial step toward this
goal rather than a mature solution for practical deployment.
